\documentclass[12pt, a4paper]{exam}
\usepackage{fontspec}
\usepackage{xunicode}
\usepackage{xltxtra}
\usepackage{polyglossia}
\usepackage{graphicx}
\usepackage{tikz}
\usetikzlibrary{positioning, calc}
\usepackage{geometry}
\usepackage{enumitem}
\usepackage{multicol}
\usepackage{array}
\usepackage{lastpage}
\usepackage{amssymb}
\usepackage{multirow}
\usepackage{amsmath}

\setmainlanguage{thai}
\XeTeXlinebreaklocale "th"
\XeTeXlinebreakskip = 0pt plus 1pt

% Page Margins
\geometry{top=0.8in, bottom=1in, left=0.8in, right=0.8in}

% Fonts
\newfontfamily\thaifont[
  Scale=1.33,
  Script=Thai,
  Path = ../fonts/,
  Extension = .ttf,
  BoldFont = *-Bold,
  ItalicFont = *-Italic,
  BoldItalicFont = *-BoldItalic
]{THSarabunNew}

\setmainfont[
  Scale=1.33,
  Path = ../fonts/,
  Extension = .ttf,
  BoldFont = *-Bold,
  ItalicFont = *-Italic,
  BoldItalicFont = *-BoldItalic
]{THSarabunNew}

% Header and Footer
\pagestyle{headandfoot}
\extraheadheight{0.2in}

\newcommand{\examheader}{%
    ชื่อ-นามสกุล \makebox[3.0in]{\hrulefill} รหัสนักศึกษา \makebox[2.0in]{\hrulefill}
}

\firstpageheader{\examheader}{}{}
\runningheader{\examheader}{}{}

\firstpagefooter{\textit{ข้อสอบรายวิชา MTH234 Linear Algebra}}{}{\thepage}
\runningfooter{\textit{ข้อสอบรายวิชา MTH234 Linear Algebra}}{}{\thepage}

% Custom Thai Choice Labels
\makeatletter
\newcommand{\thaialph}[1]{%
  \ifcase\value{#1}\or ก\or ข\or ค\or ง\or จ\else\@ctrerr\fi
}
\makeatother

% Question Format
\renewcommand{\thequestion}{\arabic{question}}
\renewcommand{\questionlabel}{\thequestion.}

% Choice Format
\renewcommand{\thechoice}{\thaialph{choice}}
\renewcommand{\choicelabel}{\thechoice)}

\renewcommand{\choiceshook}{%
  \setlength{\leftmargin}{2.5em}%
  \setlength{\labelwidth}{20pt}%
  \setlength{\labelsep}{5pt}%
}

\begin{document}
\setcounter{page}{2}

\noindent \textbf{\underline{คำสั่ง}} : ให้นักศึกษาทำข้อสอบทุกข้อ โดยทำลงไปในข้อสอบ
\vspace{0.5cm}

\begin{questions}

% Question 1: Subspaces (Abstract P3)
\question \textbf{(5 คะแนน)} กำหนดให้ $V = P_3$ เป็นปริภูมิเวกเตอร์ (vector space) ของพหุนามทั้งหมดที่มีดีกรีไม่เกิน 3 (polynomials of degree at most 3) และกำหนดให้ $W$ เป็นสับเซต (subset) ของ $V$ นิยามโดย
\[
W = \{p(x) = ax^3 + bx^2 + cx + d \in P_3 \mid 2a - b + c - 2d = 0\}
\]
จงแสดงว่า $W$ เป็นปริภูมิย่อย (subspace) ของ $P_3$
\vspace{\stretch{1}}
\newpage

% Question 2: Span, Basis, Linear Combination (R^3)
\question \textbf{(6 คะแนน)} กำหนดให้ $S = \{v_1, v_2, v_3, v_4\}$ เป็นเซตของเวกเตอร์ใน $\mathbb{R}^3$ โดยที่
\[
v_1 = \begin{bmatrix} 1 \\ 1 \\ -2 \end{bmatrix}, \quad
v_2 = \begin{bmatrix} -2 \\ -2 \\ 4 \end{bmatrix}, \quad
v_3 = \begin{bmatrix} 0 \\ 1 \\ 1 \end{bmatrix}, \quad
v_4 = \begin{bmatrix} 1 \\ 3 \\ 0 \end{bmatrix}
\]
\begin{parts}
    \part จงหาฐานหลัก (basis) สำหรับปริภูมิย่อย $W = \text{Span}(S)$ พร้อมระบุขนาดมิติ (dimension) ของ $W$
    \vspace{\stretch{2}}
    \part จงตรวจสอบว่าเวกเตอร์ $u = \begin{bmatrix} 2 \\ 5 \\ -1 \end{bmatrix}$ อยู่ใน $W$ หรือไม่ หากอยู่ จงเขียน $u$ ในรูปผลรวมเชิงเส้น (linear combination)
\end{parts}
\vspace{\stretch{1}}
\newpage

% Question 3: Solution Space (Homogeneous System)
\question \textbf{(6 คะแนน)} จงหาฐานหลัก (basis) และระบุขนาดมิติ (dimension) สำหรับปริภูมิคำตอบ (solution space) ของระบบสมการเชิงเส้นเอกพันธุ์ที่มี 4 ตัวแปรต่อไปนี้:
\[
\begin{aligned}
x_1 + x_2 + 2x_3 - 2x_4 &= 0 \\
2x_1 - x_2 - 8x_3 + 11x_4 &= 0 \\
x_1 + 2x_2 + 6x_3 - 7x_4 &= 0
\end{aligned}
\]
\vspace{\stretch{1}}
\newpage

% Question 4: Matrix Analysis with Parameter (High difficulty)
\question \textbf{(7 คะแนน)} พิจารณาเมทริกซ์ $A$ ซึ่งขึ้นอยู่กับตัวแปรเสริม $k$:
\[
A = \begin{bmatrix} 1 & 2 & k \\ 2 & 5 & 2k \\ 1 & 3 & 2 \end{bmatrix}
\]
\begin{parts}
    \part จงหาค่าแรงค์ (rank) ของ $A$ ในรูปของ $k$ (แยกกรณีตามค่า $k$)
    \vspace{\stretch{1}}
    \part จงหาค่าของ $k$ ที่ทำให้ระบบสมการ $A\mathbf{x} = \mathbf{0}$ มีผลเฉลยไม่ชัดแจ้ง (non-trivial solutions)
    \vspace{\stretch{1}}
    \newpage
    \part จากค่า $k$ ที่ได้ จงหาฐานหลักสำหรับปริภูมิสู่ศูนย์ (Null Space) และฐานหลักสำหรับปริภูมิหลัก (Column Space) ของ $A$
\end{parts}
\vspace{\stretch{1}}
\newpage

% Question 5: Change of Basis (R^3)
\question \textbf{(8 คะแนน)} กำหนดให้ $B$ และ $B'$ เป็นฐานหลัก (bases) สองฐานสำหรับปริภูมิเวกเตอร์ $\mathbb{R}^3$ โดยเขียนในรูปเมทริกซ์ (matrix form) ดังนี้:
\[
B = \left\{ \begin{bmatrix} 1 \\ 2 \\ -1 \end{bmatrix}, \begin{bmatrix} 2 \\ 3 \\ 0 \end{bmatrix}, \begin{bmatrix} 1 \\ 3 \\ -2 \end{bmatrix} \right\}
\quad \text{และ} \quad
B' = \left\{ \begin{bmatrix} 1 \\ 0 \\ 2 \end{bmatrix}, \begin{bmatrix} 2 \\ 1 \\ 3 \end{bmatrix}, \begin{bmatrix} 0 \\ 1 \\ 2 \end{bmatrix} \right\}
\]
\begin{parts}
    \part จงหาเมทริกซ์การเปลี่ยนฐานหลัก (transition matrix) $_{B'}P_{B}$ จากฐานหลัก $B$ ไปยังฐานหลัก $B'$
    \vspace{\stretch{1}}
    \newpage
    \part กำหนดให้ $v = \begin{bmatrix} 1 \\ 4 \\ -5 \end{bmatrix}$ จงหาเวกเตอร์ระบุตำแหน่ง (coordinate vector) $[v]_B$
    \vspace{\stretch{1}}
    \part จงหาเวกเตอร์ระบุตำแหน่ง (coordinate vector) $[v]_{B'}$
\end{parts}
\vspace{\stretch{1}}
\newpage

% Question 6: Orthogonality & Gram-Schmidt (R^4)
\question \textbf{(8 คะแนน)} กำหนดให้ $W$ เป็นปริภูมิย่อยของ $\mathbb{R}^4$ สแปนโดยเวกเตอร์ (spanned by the vectors)
\[
w_1 = \begin{bmatrix} 1 \\ 2 \\ 1 \\ 1 \end{bmatrix}, \quad
w_2 = \begin{bmatrix} 2 \\ 3 \\ 0 \\ 2 \end{bmatrix}, \quad
w_3 = \begin{bmatrix} 4 \\ 2 \\ 2 \\ 0 \end{bmatrix}
\]
และให้ $v = \begin{bmatrix} -1 \\ 1 \\ 2 \\ 3 \end{bmatrix}$
\begin{parts}
    \part จงใช้กระบวนการ Gram-Schmidt (Gram-Schmidt process) เพื่อหาฐานหลักเชิงตั้งฉากปกติ (orthonormal basis) สำหรับ $W$
    \vspace{\stretch{1}}
    \newpage
    \part จงหาการฉายเชิงตั้งฉาก (orthogonal projection) ของ $v$ บน $W$ ($\text{proj}_W v$)
\end{parts}
\vspace{\stretch{1}}

\end{questions}
\end{document}
